% Working Notes on Chiral Koszul Duality
% Raeez Lorgat, 2024--2026
%
% Compile: pdflatex working_notes.tex (standalone)

\documentclass[11pt]{article}

\usepackage[T1]{fontenc}
\usepackage[utf8]{inputenc}
\usepackage{lmodern}
\frenchspacing

\usepackage[cmintegrals,cmbraces,noamssymbols]{newtxmath}
\usepackage{ebgaramond}

\usepackage[
  activate={true,nocompatibility},
  final,
  tracking=true,
  kerning=true,
  spacing=true,
  factor=1100,
  stretch=10,
  shrink=10
]{microtype}

\usepackage{mleftright}
\mleftright

\let\Bbbk\undefined
\let\openbox\undefined
\usepackage{amsmath,amssymb,amsthm}
\usepackage{mathtools}
\usepackage{enumitem}
\usepackage[margin=1.2in]{geometry}
\usepackage[unicode]{hyperref}
\usepackage{tikz-cd}

\setlength{\emergencystretch}{3em}
\raggedbottom

% Theorem styles
\usepackage{thmtools}
\declaretheoremstyle[
  spaceabove=\topsep, spacebelow=\topsep,
  headfont=\normalfont\scshape,
  notefont=\normalfont\itshape,
  bodyfont=\normalfont,
  postheadspace=0.5em, headpunct={.}
]{notesthm}
\declaretheoremstyle[
  spaceabove=\topsep, spacebelow=\topsep,
  headfont=\normalfont\itshape,
  notefont=\normalfont\itshape,
  bodyfont=\normalfont,
  postheadspace=0.5em, headpunct={.}
]{notesdef}

\declaretheorem[style=notesthm, name=Theorem, numberwithin=section]{theorem}
\declaretheorem[style=notesthm, name=Lemma, sibling=theorem]{lemma}
\declaretheorem[style=notesthm, name=Proposition, sibling=theorem]{proposition}
\declaretheorem[style=notesthm, name=Corollary, sibling=theorem]{corollary}
\declaretheorem[style=notesdef, name=Definition, sibling=theorem]{definition}
\declaretheorem[style=notesdef, name=Example, sibling=theorem]{example}
\declaretheorem[style=notesdef, name=Remark, sibling=theorem]{remark}
\declaretheorem[style=notesdef, name=Observation, sibling=theorem]{observation}
\declaretheorem[style=notesdef, name=Warning, sibling=theorem]{warning}
\declaretheorem[style=notesdef, name=Conjecture, sibling=theorem]{conjecture}
\declaretheorem[style=notesdef, name=Question, sibling=theorem]{question}

% Macros
\newcommand{\cA}{\mathcal{A}}
\newcommand{\cB}{\mathcal{B}}
\newcommand{\cM}{\mathcal{M}}
\newcommand{\cO}{\mathcal{O}}
\newcommand{\cW}{\mathcal{W}}
\newcommand{\cH}{\mathcal{H}}
\newcommand{\cZ}{\mathcal{Z}}
\newcommand{\cD}{\mathcal{D}}
\newcommand{\fg}{\mathfrak{g}}
\newcommand{\fh}{\mathfrak{h}}
\newcommand{\fn}{\mathfrak{n}}
\newcommand{\barB}{\overline{B}}
\newcommand{\B}{\bar{B}}
\newcommand{\C}{\mathbb{C}}
\newcommand{\Z}{\mathbb{Z}}
\newcommand{\Q}{\mathbb{Q}}
\newcommand{\bP}{\mathbb{P}}

\newcommand{\Ass}{\mathsf{Ass}}
\newcommand{\Com}{\mathsf{Com}}
\newcommand{\Lie}{\mathsf{Lie}}
\newcommand{\Eone}{\mathsf{E}_1}
\newcommand{\Etwo}{\mathsf{E}_2}
\newcommand{\Einf}{\mathsf{E}_{\infty}}
\newcommand{\Linf}{\mathsf{L}_{\infty}}
\newcommand{\Ainf}{\mathsf{A}_{\infty}}
\newcommand{\Pinf}{\mathsf{P}_{\infty}}
\newcommand{\chirAss}{\mathsf{Ass}^{\mathrm{ch}}}
\newcommand{\chirCom}{\mathsf{Com}^{\mathrm{ch}}}
\newcommand{\chirLie}{\mathsf{Lie}^{\mathrm{ch}}}
\newcommand{\Ran}{\operatorname{Ran}}
\newcommand{\Res}{\operatorname{Res}}
\newcommand{\FM}{\mathrm{FM}}
\newcommand{\CE}{\mathrm{CE}}
\newcommand{\DS}{\mathrm{DS}}
\newcommand{\MC}{\mathrm{MC}}
\newcommand{\ad}{\operatorname{ad}}
\newcommand{\id}{\mathrm{id}}
\newcommand{\Hom}{\operatorname{Hom}}
\newcommand{\Ext}{\operatorname{Ext}}
\newcommand{\Defcyc}{\operatorname{Def}_{\mathrm{cyc}}}
\newcommand{\dzero}{d_0}
\newcommand{\dfib}{d_{\mathrm{fib}}}
\newcommand{\Dg}[1]{D_{#1}}
\newcommand{\Dtot}{D_{\mathrm{tot}}}
\newcommand{\normord}[1]{:\!#1\!:}
\newcommand{\Fact}{\mathsf{Fact}}

\DeclareMathOperator{\gr}{gr}
\DeclareMathOperator{\Bl}{Bl}
\DeclareMathOperator{\Spec}{Spec}

\title{\textit{Working Notes on Chiral Koszul Duality}\\[0.4em]
\large Dead Ends, Faery Gates, and Nuclear Warheads}
\author{Raeez Lorgat}
\date{2024--2026}

\begin{document}
\maketitle

\begin{quote}
\small\itshape
A differential that refuses to square to zero.  A sign audit
across 2048 conventions that finds nothing.  A genus-one propagator
that bends the Arnold relation by exactly one $(1,1)$-form.
A finite-dimensional thick generation argument that dies on
contact with an infinite-dimensional Verma module.  A cubic Casimir
that everyone expected to contribute at genus two but sits
in the wrong degree of the cyclic complex.

These are the moments when the subject taught us what it
actually is.  This document records them.
\end{quote}

\tableofcontents

% =================================================================
\section{On a differential that does not square to zero}
\label{sec:the-crime}
% =================================================================

\noindent
The crime scene is the bar complex of an affine Kac--Moody algebra.

\subsection{The scene}

A chiral algebra $\cA$ on a smooth curve $X$ has an operator product
expansion with poles.  The Heisenberg current has a double pole:
$\alpha(z)\alpha(w) = k/(z-w)^2$.  The Kac--Moody current $J^a(z)$
has both:
\begin{equation}\label{eq:km-ope}
J^a(z) J^b(w) = \frac{k\kappa^{ab}}{(z-w)^2}
  + \frac{f^{ab}{}_c\, J^c(w)}{z-w} + \text{regular}.
\end{equation}
The simple pole is the Lie bracket.  The double pole is the
invariant form---the ``level.''  These are different organs of
the same creature.

The bar construction places $n$ elements of $\cA$ at $n$ points of
$X$, dresses them with logarithmic $1$-forms
$\eta_{ij} = d\log(z_i - z_j)$, and extracts residues at
collisions.  Out falls a differential $d$ on the graded space
$\B(\cA) = \bigoplus_n (s^{-1}\bar{\cA})^{\otimes n} \otimes
\Omega^{n-1}_{\log}(\overline{C}_n(X))$.

Decompose $d$ by pole order:
\begin{align}
d_{\mathrm{bracket}} &\colon
  \text{extracts the simple pole } (\phi_i)_{(0)}\phi_j,
  \label{eq:d-bracket}\\
d_{\mathrm{curvature}} &\colon
  \text{extracts the double pole } (\phi_i)_{(1)}\phi_j.
  \label{eq:d-curv}
\end{align}

In classical homological algebra, the bracket differential of a
Lie algebra $\fg$ squares to zero by the Jacobi identity.  Every
algebraist's spinal cord expects $d_{\mathrm{bracket}}^2 = 0$.

\begin{theorem}[The crime]\label{thm:the-crime}
$d_{\mathrm{bracket}}^2 \neq 0$ for any chiral algebra with
both a simple and a double pole in the OPE.  This was verified
across all $2048$ sign conventions arising from eleven
independent binary choices \textup{(}grading, suspension,
associativity, antisymmetry, contraction, composition,
Koszul signs\textup{)}.  Every convention fails.
\end{theorem}

This is not a sign error.  It is a theorem.

\subsection{The mechanism}

When $d_{\mathrm{bracket}}$ is applied twice, the composition
passes through iterated residues at codimension-$2$ strata---the
loci where three points collide.  The Jacobi identity cancels the
pure bracket--bracket terms, but there are \emph{bracket--curvature
cross-terms}: the Taylor expansion of the OPE near the triple
collision locus $D_{ijk}$ mixes pole orders.  Concretely, the
iterated residue
$\Res_{D_{j\ell}} \Res_{D_{ij}}
[(\phi_i)_{(0)}((\phi_j)_{(1)}\phi_\ell \cdot |0\rangle)]$
does not vanish, and no choice of signs can make it vanish, because
it is a nonzero number (proportional to the level $k$) multiplied by
a nonzero vector.

The identity that \emph{does} hold:
\begin{equation}\label{eq:borcherds}
d_{\mathrm{res}}^2 = (d_{\mathrm{bracket}} + d_{\mathrm{curvature}})^2
= 0.
\end{equation}
Since $d_{\mathrm{curvature}}^2 = 0$ (the double-pole output is a
scalar times the vacuum; a subsequent extraction sees
$|0\rangle_{(n)}\phi = 0$), we get:
\[
d_{\mathrm{bracket}}^2
= -\{d_{\mathrm{bracket}}, d_{\mathrm{curvature}}\}.
\]
The bracket differential fails to square to zero by
\emph{exactly the anticommutator with the curvature differential}.
The pair $(d_{\mathrm{bracket}}, d_{\mathrm{curvature}})$ is not
a bicomplex.  The level is the obstruction to
$d_{\mathrm{bracket}}$ being a differential.

\begin{remark}[The physics]
This identity \emph{is} the conformal anomaly.  The BRST operator
alone ($d_{\mathrm{bracket}}$) does not square to zero; the anomaly
term ($d_{\mathrm{curvature}}$) is what prevents it; the \emph{full}
BRST operator ($d_{\mathrm{res}}$) squares to zero because it
includes both.  The quantum master equation
$\hbar\Delta S + \tfrac{1}{2}\{S,S\} = 0$ is the same identity
written in BV language.  The ``obstruction'' is the content.
\end{remark}

\subsection{Why geometry saves what algebra cannot}

\begin{observation}[Arnold is the gatekeeper]
The full differential $d_{\mathrm{res}}$ squares to zero because
the geometry of the Fulton--MacPherson compactification
$\overline{C}_n(X)$ enforces it.  The computation of $d^2$ produces
triple residues at codimension-$2$ strata $D_{ijk}$, and the
\emph{Arnold relation}
\[
\eta_{12} \wedge \eta_{23}
+ \eta_{23} \wedge \eta_{31}
+ \eta_{31} \wedge \eta_{12} = 0
\]
makes these vanish.  This is Stokes' theorem on a manifold with
corners: at the corner $D_{ij} \cap D_{jk}$, the two adjacent
faces contribute opposite orientation signs.  The cancellation is
topological; no algebraic sign convention can reproduce it or
circumvent it.

The Borcherds identity for vertex algebras is the algebraic shadow
of this geometric cancellation.  It governs the full OPE
simultaneously---all poles at once---because Stokes' theorem does
not care about pole order.
\end{observation}

\begin{remark}[The propagator is unique]
The form $\eta_{ij} = d\log(z_i - z_j)$ is the \emph{unique}
meromorphic $1$-form on $C_2(X)$ satisfying three simultaneous
constraints: conformal invariance (unchanged under $z \mapsto f(z)$),
well-defined residues (a simple pole has a coordinate-independent
residue; a double pole does not), and the Arnold relation (ensuring
$d^2 = 0$).  No other $1$-form satisfies all three.  The propagator
is not a choice.  It is the unique gate through which the OPE data of
a chiral algebra can pass into a cochain complex.
\end{remark}


% =================================================================
\section{Three differentials, three worlds}
\label{sec:three-worlds}
% =================================================================

\noindent
For eighteen months, the letter $d$ was used for three different
operators.  This notational sin hid the deepest structural feature
of the subject.

\subsection{The trinity}

On a genus-$g$ bar complex, three differentials act:

\medskip
\noindent
\textbf{(i) The genus-$0$ collision differential $\dzero$.}
Built from residues on $\bP^1$ using
$\eta_{ij} = d\log(z_i - z_j)$.  Satisfies $\dzero^2 = 0$ by
Arnold.  This is the differential of classical Koszul duality.

\medskip
\noindent
\textbf{(ii) The fiberwise curved differential $\dfib$.}
On a fixed curve $\Sigma_g$, the propagator acquires period
corrections:
\[
\eta^{(g)}_{ij} = d\log E(z_i, z_j)
  + \sum_{k=1}^g \omega_k(z_i) \oint_{A_k} \omega_k(z_j),
\]
where $E(z_i, z_j)$ is the prime form (a section of
$K^{-1/2} \boxtimes K^{-1/2}$, not $K^{+1/2}$---this sign costs
a week if wrong) and $\omega_k$ is a basis of
$H^0(\Sigma_g, \Omega^1)$.

The Arnold relation \emph{fails}.  The residual is the Arakelov
$(1,1)$-form:
\[
\dfib^{\,2} = \kappa(\cA) \cdot \omega_g.
\]
The bar differential no longer squares to zero.  The world is curved.

\medskip
\noindent
\textbf{(iii) The total corrected differential $\Dg{g}$.}
$\Dg{g} = \dzero + \sum_{k=1}^g t_k\, d_k$,
where $t_k = \oint_{A_k}\omega_k$ are $A$-period parameters from
$H^1(\Sigma_g, \C)$.
Satisfies $\Dg{g}^2 = 0$.

\begin{theorem}[Lagrangian cancellation]
\label{thm:lagrangian}
The $A$-cycle subspace of $H^1(\Sigma_g, \C)$ is Lagrangian for the
intersection pairing.  The curvature $\dfib^2$ lies in the
symplectic complement.  The $A$-period corrections project it to
zero.  This is why $\Dg{g}^2 = 0$: the correction terms and the
curvature annihilate each other because they live in complementary
Lagrangian subspaces.
\end{theorem}

\begin{warning}[The parameter-source confusion]
\label{warn:not-moduli}
The period parameters $t_k$ live in $H^1(\Sigma_g, \C)$---periods
of the \emph{fixed curve}.  They are \textbf{not} classes on
$\cM_g$.  This matters because $H^1(\cM_g, \Q) = 0$ for
$g \geq 2$ (Harer).  There are no such classes.

The scalar obstruction $\kappa(\cA) \cdot \lambda_g$ lives in the
tautological ring of $\overline{\cM}_g$, at the end of a chain:
\[
H^1(\Sigma_g, \C)
\xrightarrow{\;\MC\;}
Z^1(\Defcyc(\cA))
\xrightarrow{\;\text{global}\;}
R\Gamma(\overline{\cM}_g, \cZ_\cA)
\xrightarrow{\;\mathrm{tr}\;}
R\Gamma(\overline{\cM}_g, \Q).
\]
Confusing fiber with base is the most common error in the subject.
Three early drafts contained it.
\end{warning}

\subsection{The curvature is the content}

Physicists read $\dfib^2 = \kappa \cdot \omega_g$ as a failure:
``the theory is anomalous.''  But the anomaly \emph{is} the modular
characteristic of the chiral algebra.  The number $\kappa(\cA)$
controls all genus-$g$ phenomena.

In the curved $\Ainf$ framework, the curvature element is
$m_0 = \kappa(\cA) \cdot \omega_g$, and:
\[
m_1^2(a) = m_2(m_0, a) - m_2(a, m_0) = [m_0, a]
\qquad\text{(commutator, minus sign)}.
\]
The ``failure'' of $m_1^2 = 0$ is the assertion that the vacuum has
a nontrivial bracket with everything.  This is the algebraic encoding
of the anomaly.  The condition $\kappa = 0$ is simultaneously
$\dfib^2 = 0$ (uncurvedness) and $c = 26$ (anomaly cancellation for
the bosonic string).

The generating function of the obstruction classes
$\kappa(\cA) \cdot \lambda_g$ across all genera is
$\kappa(\cA) \cdot (\hat{A}(i\hbar) - 1)$: the $\hat{A}$-genus of
an abstract Dirac operator whose index is the modular characteristic.
This is the family index theorem of Bismut--Freed, specialized to
the Hodge bundle over $\overline{\cM}_g$.  Its appearance from OPE
residues is not visible from the definition.  The bar construction
knows things you did not tell it.


% =================================================================
\section{On the faery structure of $\Eone$-chiral algebras}
\label{sec:e1-chiral}
% =================================================================

\noindent
There are two axes called ``$\Eone$,'' and confusing them will
destroy you.

\subsection{Two axes, one notation, total darkness}

\medskip
\noindent
\textbf{Axis 1: The chiral hierarchy} (fixed curve, varying
commutativity).

On a smooth curve $X$, three species of chiral algebra coexist:
\[
\Einf\text{-chiral}
\;\subset\; \Pinf\text{-chiral}
\;\subset\; \Eone\text{-chiral}.
\]

An $\Einf$-chiral algebra is a vertex algebra in the sense of
Beilinson--Drinfeld: a $\cD_X$-module with a skew-symmetric chiral
bracket
$\mu\colon j_*j^*(\cA \boxtimes \cA) \to \Delta_!\cA$.
The skew-symmetry axiom $Y(a,z)b = e^{z\partial}Y(b,-z)a$ forces
full locality.  Configuration spaces are \emph{unordered}.

An $\Eone$-chiral algebra drops skew-symmetry.  The chiral bracket
is \emph{associative} but not commutative.  Fields satisfy
$R$-twisted locality:
\begin{equation}\label{eq:r-twisted}
J_a(z) J_b(w) - R^{ab}_{cd}(z-w)\, J_d(w) J_c(z) = 0,
\end{equation}
where $R(u)$ is a matrix-valued rational function---the Yang
$R$-matrix.  When $R = P$ (the permutation), this is
skew-symmetry; we recover $\Einf$.  When $R \neq P$, it is
strictly $\Eone$: the fields do not commute, and configuration
spaces must be \emph{ordered}.

A $\Pinf$-chiral algebra has both an $\Einf$ product and a
$\chirLie$ bracket satisfying the Leibniz rule.  It is the chiral
Poisson algebra, intermediate between the two extremes.

\medskip
\noindent
\textbf{Axis 2: The topological ladder} (varying manifold dimension,
$\Eone/\Etwo/\ldots/\mathsf{E}_n$).

The little $n$-disks operad $\mathsf{E}_n$ lives on real
$n$-manifolds.  Arnold relations (three pairwise collisions,
$n = 2$) generalize to Totaro relations (higher codimension,
$n \geq 3$).  A complex curve has complex dimension~$1$ but
\emph{real} dimension~$2$: it enters the topological ladder at
$n = 2$, not $n = 1$.

\begin{warning}[The fatal confusion]
\label{warn:e1-axes}
An $\Eone$-\emph{chiral} algebra lives on a curve (complex
dimension~$1$, real dimension~$2$).  It is \textbf{not} a
topological $\Eone$-algebra (which lives on an interval, real
dimension~$1$).  The suffix ``$\text{-chiral}$'' is
mathematically load-bearing.  Without it, the notation is a
death trap.

The chiral bar complex on a curve is an $\Etwo$-algebraic
object (Arnold relations from the real $2$-dimensional configuration
space $C_n(\Sigma_g)$) enriched by holomorphic structure (the OPE
poles are holomorphic, not just topological).  It sits at the
$\Etwo$ level of the topological ladder with a holomorphic
refinement.

At the topological $\Eone$ level (intervals/circles), the bar
complex reduces to the classical $\Ainf$ bar-cobar adjunction:
no configuration space geometry, no Arnold relation, no propagator.
The entire subject of this monograph vanishes.
\end{warning}

\subsection{Why ordered configurations are forced: a fairy tale in four acts}

Here is a story.  It has the logic of a fairy tale and the payload
of a nuclear weapon.

\medskip
\noindent
\textsc{Act I.}
The Kazhdan--Lusztig equivalence sends integrable
$\widehat{\fg}_k$-modules to representations of the quantum group
$U_q(\fg)$, where $q = e^{\pi i/(k+h^\vee)}$.  The quantum group
has a braiding $\sigma$ given by the universal $R$-matrix.
\emph{This braiding is not symmetric}: $R_{12} \neq P_{12}$.

\medskip
\noindent
\textsc{Act II.}
Bar-cobar duality on the Kac--Moody side acts by Verdier duality on
$\overline{C}_n(X)$.  On the module category, this \emph{reverses the
braiding}: $\sigma \mapsto \sigma^{-1}$.  Under the KL equivalence,
the reversed braiding corresponds to $R \mapsto R^{-1}$, which is
$q \mapsto q^{-1}$.

\medskip
\noindent
\textsc{Act III.}
On the quantum group side, $R$-matrix inversion is not level
inversion.  It is a Hopf algebra involution that reverses the ordering
of tensor factors.  To realize this geometrically on the curve, we
must track which insertion comes first: left or right.  This forces
\emph{ordered} configuration spaces $C_n^{\mathrm{ord}}(X)$.

\medskip
\noindent
\textsc{Act IV.}
Without ordering, Verdier duality acts by level inversion
$k \mapsto -k - 2h^\vee$ (the $\Einf$ story, recovering
Feigin--Frenkel).  \emph{With} ordering, it acts by $R$-matrix
inversion $R \mapsto R^{-1}$ (the $\Eone$ story, recovering the
Drinfeld--Kohno theorem).  The Yangian $Y(\fg)$, presented via the
RTT relation
$R_{12}(u-v)\, T_1(u)\, T_2(v) = T_2(v)\, T_1(u)\, R_{12}(u-v)$,
is the archetypal $\Eone$-chiral algebra: its Koszul dual is
$Y_{R^{-1}}(\fg)$, the same Yangian with inverted $R$-matrix.

\begin{theorem}[$\Eone$-chiral Koszul duality]
\label{thm:e1-koszul}
For the Yangian $Y_\hbar(\mathfrak{sl}_N)$ in RTT presentation:
\[
Y_\hbar(\mathfrak{sl}_N)^{!,\mathrm{ch}}
\cong Y_{R^{-1},\hbar}(\mathfrak{sl}_N)^{\mathrm{ch}},
\]
where the dual RTT relation uses the inverse Yang $R$-matrix:
$R^{-1}(u) = 1 + \hbar P/u + O(u^{-2})$ \textup{(}sign flip in the
linear term\textup{)}.
\end{theorem}

This is a nuclear warhead wrapped in a fairy tale.  The passage from
$R$ to $R^{-1}$ is the $\Eone$-chiral Koszul duality, and it
unifies:
\begin{itemize}[nosep]
\item the Feigin--Frenkel involution (the $\Einf$ shadow),
\item the Drinfeld--Kohno equivalence between KZ monodromy and
  quantum group $R$-matrices,
\item the classical $\Ass^! \cong \Ass$ self-duality of the
  associative operad (the point-limit shadow).
\end{itemize}

\subsection{The cocycle gate: $\Einf$ to $\Eone$ in one parameter}

Lattice vertex algebras make the $\Einf \leftrightarrow \Eone$
transition visible in a single parameter.  A $2$-cocycle
$\varepsilon$ on an even lattice $\Lambda$ satisfying the Borcherds
symmetry axiom
$\varepsilon(\alpha,\beta) = (-1)^{\langle\alpha,\beta\rangle}
\varepsilon(\beta,\alpha)$ produces a vertex algebra
$V_\Lambda^\varepsilon$.

When $\varepsilon$ is \emph{symmetric}:
$V_\Lambda^\varepsilon$ is $\Einf$-chiral.
When $\varepsilon$ is \emph{not symmetric}:
$V_\Lambda^\varepsilon$ is strictly $\Eone$-chiral.

An irrational twist $\varepsilon((a,b),(c,d)) = e^{2\pi i\theta \cdot ad}$
on $\Z^2$ produces the noncommutative torus in chiral dress: a
single irrational number $\theta$ turns a commutative algebra into a
nonlocal one.  Koszul duality inverts it:
$\varepsilon \mapsto \varepsilon^{-1}$.

The $\Eone$ inversion principle---$\varepsilon \mapsto
\varepsilon^{-1}$ for lattices, $R \mapsto R^{-1}$ for
Yangians---is the same phenomenon at different levels of
generality.  The lattice case is the abelian shadow of the Yangian
case, just as the Heisenberg is the abelian shadow of Kac--Moody.

\subsection{Coisson is not chiral Poisson}

A confusion that has ruined entire papers:

\begin{itemize}[nosep]
\item A \textbf{Coisson algebra} (Beilinson--Drinfeld) $=$ a Poisson
  vertex algebra $=$ a commutative $\cD_X$-algebra with
  $\lambda$-bracket.  It is \textbf{not} a chiral algebra.  It lives
  outside the chiral hierarchy.
\item A \textbf{$\Pinf$-chiral algebra} $=$ $\chirCom + \chirLie$ with
  Leibniz.  It \textbf{is} a chiral algebra---in fact an
  $\Einf$-chiral algebra equipped with additional Lie structure.
\end{itemize}

The quantization levels are different:
\[
\text{Coisson}
\;\xrightarrow{\text{singly quantum}}\;
\Einf\text{-chiral}
\;\xrightarrow{\text{doubly quantum}}\;
\Eone\text{-chiral}.
\]
Coisson $\to$ $\Einf$ is the passage from Poisson to associative
(star product).  $\Einf \to \Eone$ is the passage from commutative
to noncommutative ($R$-matrix deformation).  These are consecutive
floors, not the same floor.  Khan--Zeng's 3d holomorphic-topological
gauge theories produce Poisson vertex algebras (Coisson); their
quantization lands at the $\Einf$ level, not the $\Eone$ level.
To reach $\Eone$ requires a further deformation: the $R$-matrix.


% =================================================================
\section{Holomorphic-topological descent and the 3d--2d bridge}
\label{sec:ht}
% =================================================================

\noindent
The relationship between $3$-dimensional theories and $2$-dimensional
chiral algebras is a bridge with specific load-bearing capacities.

\subsection{Costello--Li: 4d to 2d via holomorphic-topological twist}

Start with $4$d $\mathcal{N}=2$ super Yang--Mills on $\C^2$.  After
the holomorphic-topological twist, the fields become
$\bar\partial$-closed forms valued in $\fg \otimes \cO_{\C^2}$, and
the action becomes $S = \{Q_{\mathrm{BRST}}, \Psi\} + S_{\mathrm{top}}$.
Replace $\C^2$ by $\C \times \C^*$, take $S^1$-equivariant reduction
along $|w| = 1$.  Out falls a \emph{factorization algebra} on $\C$.

Is it a chiral algebra?  Beilinson--Drinfeld says yes, under the
equivalence between factorization algebras and chiral algebras on
curves.  But this equivalence requires the factorization algebra to
have prescribed pole structure along the diagonal, which is not
automatic from the reduction.  The conditions:
holomorphicity of the Coulomb branch, survival of conformal symmetry,
and $\cD$-module compatibility.

When these conditions hold, the Costello--Li reduction produces a
genuine chiral algebra in the sense of BD, and the entire bar-cobar
machinery applies.  When they fail, you have a factorization algebra
that is not a chiral algebra, and the bar complex must be built in
the factorization framework (Costello--Gwilliam) rather than the
$\cD$-module framework (BD/FG/this monograph).

\subsection{Paquette--Williams: boundary vertex algebras}

The $4$d holomorphic-topological theory on $\C^2$ with a boundary
at $z_2 = 0$ supports a \emph{vertex operator algebra} $V_{\cB}$
on the boundary.  The boundary VOA has a \emph{chiral envelope}
$\cA_{\cB}$---a genuine chiral algebra in the BD sense.

The bar-cobar framework applies to the chiral envelope, not to the
VOA directly (the VOA is an algebraic object; the chiral envelope
is a geometric $\cD$-module object).  The distinction matters: the
bar complex of $\cA_{\cB}$ computes factorization homology on
$\Ran(X)$, which is the correct ``spectrum'' of the boundary theory.

\subsection{The 3d holomorphic-topological zoo}

Three dimensional holomorphic-topological theories produce structure
on the first floor of the quantization elevator:

\begin{center}
\small
\begin{tabular}{lll}
\textbf{Origin} & \textbf{Structure} & \textbf{Chiral level} \\
\hline
3d HT gauge (Khan--Zeng) & Poisson vertex algebra &
  Coisson (pre-chiral) \\
3d $\mathcal{N}=4$ boundary & VOA on boundary &
  $\Einf$-chiral envelope \\
3d Chern--Simons & Config.\ space integrals ($n=3$) &
  $\mathsf{E}_3$-algebraic \\
4d HT $\to$ 2d reduction & Factorization algebra on $\C$ &
  $\Einf$-chiral if conditions hold
\end{tabular}
\end{center}

\noindent
The $\Eone$-chiral level---Yangians, quantum groups, non-symmetric
$R$-matrices---does not arise directly from $3$d HT gauge theory.
It requires the additional ``doubly quantum'' deformation of the
$R$-matrix.  This is the Drinfeld--Jimbo quantization of the
classical $r$-matrix, which is a $2$-dimensional phenomenon (KZ
equations on a curve), not a $3$-dimensional one.

\begin{remark}[Chern--Simons and the Arnold--Totaro transition]
In $3$d Chern--Simons theory, the relevant configuration spaces
are $C_n(M^3)$ for a $3$-manifold $M$, and the structural relations
are \emph{Totaro relations} (generalizing Arnold from $n = 2$ to
$n = 3$).  The bar complex on $C_n(M^3)$ computes Vassiliev
invariants via BV integration.  The bridge to $2$d chiral algebras
comes through the holomorphic Chern--Simons/WZW correspondence:
the CS path integral on $\Sigma_g \times [0,1]$ produces the
WZW chiral algebra on $\Sigma_g$.

The perturbative CS expansion gives $\Ainf$ structure on the
deformation complex of flat connections.  This is the $3$d avatar of
the curved $\Ainf$ structure on the genus-$g$ bar complex.  The
curvature $m_0$ in $2$d corresponds to the CS level in $3$d; the
condition $\kappa = 0$ (anomaly cancellation) corresponds to the
integrality of the CS level.
\end{remark}


% =================================================================
\section{Chiral Koszulness is not classical Koszulness}
\label{sec:koszulness}
% =================================================================

\noindent
$U(\fg)$ is Koszul.  $\widehat{\fg}_k$ might not be.

\subsection{The distinction}

Classical Koszulness: the bar complex $\B(U(\fg))$ has cohomology
concentrated in a single row.  This uses the associative product
(a single binary operation).  Proved by Priddy, Backelin--Fr\"oberg,
BGS.

Chiral Koszulness: the \emph{chiral} bar complex
$\B^{\mathrm{ch}}(\widehat{\fg}_k)$, built from the full OPE
including all poles, has the analogous concentration.  This uses
all pole orders simultaneously (an infinite family of operations).

Classical does not imply chiral.  The higher OPE poles are not a
perturbation; they are structural.  The Sugawara tensor
$T = \frac{1}{2(k+h^\vee)}\sum_a \normord{J^a J^a}$ involves
double Wick contractions producing quartic poles in $T(z)T(w)$, and
these generate new differentials in the PBW spectral sequence with
no classical analogue.

\subsection{The circularity}

To compute bar cohomology dimensions, assume chiral Koszulness (the
spectral sequence collapses).  But the bar cohomology dimensions
are what would verify chiral Koszulness (via
$H_\cA(t) \cdot H_{\cA^!}(-t) = 1$).  The circle must be broken
by an \emph{independent} collapse proof---family-specific, using:
\begin{itemize}[nosep]
\item \textbf{Kac--Moody}: Whitehead's theorem + Killing contraction.
  The enrichment is $\fg \otimes \fg$ (dimension $\dim(\fg)^2$),
  \textbf{not} $\Lambda^2(\fg)$
  (dimension $\binom{\dim\fg}{2}$).  Getting this wrong
  produces systematically incorrect predictions.  Several months
  were lost to this error.
\item \textbf{Virasoro}: The $\mathfrak{sl}_2$ subalgebra
  $\{L_{-1}, L_0, L_1\}$ acts $c$-independently (the central term
  $m^3 - m = 0$ for $m \in \{-1,0,1\}$).  Casimir eigenvalues at
  weight $h$ are all positive for $h \geq 2$.
\item \textbf{Principal $\cW_N$}: Unique weight-$2$ mechanism from
  the stress tensor, triangular $d_2$ argument.
\end{itemize}

There is no universal collapse theorem.  The three mechanisms share
a family resemblance but not a common proof.

\subsection{What free fields teach}

For free fields (Heisenberg, free fermion), the OPE has only a
double pole: $d_{\mathrm{bracket}} = 0$, so $d = d_{\mathrm{curvature}}$,
which squares to zero automatically.  No Borcherds identity needed.
No PBW spectral sequence needed.  No circularity.

The jump from ``one pole'' (free) to ``two poles'' (interacting) is
not incremental.  It is a phase transition.  The entire machinery of
the monograph---the Borcherds identity, the PBW spectral sequence,
the Arnold--Stokes--Verdier triangle, the family-specific collapse
proofs---exists to handle this transition.  Free-field computations
are a warm-up in a different sport.


% =================================================================
\section{The Heisenberg frame atom}
\label{sec:heisenberg-frame}
% =================================================================

\noindent
Every theorem of the monograph is visible in the Heisenberg algebra.
None of the pitfalls are.  This is why it is the frame atom.

The OPE $\alpha(z)\alpha(w) = k/(z-w)^2$ has a double pole and
nothing else.  The bar differential is pure curvature.  The bar
cohomology at genus $0$ counts partitions.  The Koszul dual is
$\cH_k^! = \mathrm{Sym}^{\mathrm{ch}}(V^*)$, \textbf{not}
$\cH_{-k}$: the Heisenberg is not self-dual.  (For Kac--Moody,
$\widehat{\fg}_k^! \cong \widehat{\fg}_{-k-2h^\vee}$ \emph{looks}
like a level flip, but this is specific to the KM OPE structure,
not a general principle.)

The genus-$1$ curvature is $\dfib^2 = k \cdot \omega_1$, and the
Fourier--Mukai identification (following Polishchuk) makes the bar
complex on $E_\tau$ literally the Fourier--Mukai kernel on
$\mathrm{Jac}(E_\tau) \times \mathrm{Jac}(E_\tau)^\vee$.  The
chiral bar construction \emph{is} the Fourier--Mukai transform,
generalized from abelian (Heisenberg) to non-abelian (Kac--Moody,
Virasoro, Yangian).

Complementarity:
$c(\cH_k) + c(\cH_k^!) = 1 + 1 = 2 = 2\dim(V)$.  For Kac--Moody:
$c + c' = 2\dim(\fg)$.  The sum is topological.

The $\hat{A}$-genus appears as the generating function of
$\kappa(\cA) \cdot \lambda_g$ across genera.  It is the
characteristic function of the Fourier transform---the analogue of
$e^{-t|\xi|^2}$ for the classical heat kernel.


% =================================================================
\section{Scalar saturation, or: the cubic Casimir that wasn't there}
\label{sec:saturation}
% =================================================================

\noindent
The Maurer--Cartan element $\Theta_\cA$ was expected to be
complicated.  It is not.

For \emph{any} simple Lie algebra $\fg$, regardless of rank:
\[
\Theta^{\min}_\cA = \kappa(\cA) \cdot \eta \otimes \Lambda.
\]
One number $\kappa(\cA)$, one universal kernel, one tautological
class.  The higher Casimirs do not contribute.

The reason: $H^2_{\mathrm{cyc}}(\fg, \fg) = H^3(\fg) = \C$
is one-dimensional for \emph{every} simple $\fg$.  The cubic
Casimir of $\mathfrak{sl}_3$ lives in $H^5(\fg)$ (degree~$5$ in
cyclic cohomology)---the wrong degree for the MC equation, which
lives at degree~$2$.

Earlier versions of the manuscript contained a proposition
claiming scalar saturation holds only for $\mathfrak{sl}_2$, and a
remark claiming the cubic Casimir would contribute non-scalar
corrections at higher rank.  Both were wrong.  The degree mismatch
is absolute: higher Casimirs sit in the wrong cohomological degree
\emph{regardless} of how large the rank is.

\begin{theorem}[Genus rigidity]
Genera $1$ and $2$ are rigid: no freedom beyond $\kappa(\cA)$.
The first nontrivial correction appears at genus~$3$:
\[
l_1(\theta_3 - \kappa(\cA)\mu \otimes \lambda_3)
= -\tfrac{1}{6}\, l_3(\theta_1, \theta_1, \theta_1),
\]
and this is purely in the \emph{Killing $l_3^\eta$ channel}---the
cubic $\Linf$ operation built from the Killing $3$-cocycle.
\end{theorem}

The non-scalar data in the modular package lives in $\Delta_\cA$
(the spectral discriminant) and $\Pi_\cA$ (the periodicity operator),
not in $\Theta_\cA$.  The universal class is scalar.  The world is
simpler than we feared.


% =================================================================
\section{The Category $\cO$ generation problem, \\
or: a theorem that died on arrival}
\label{sec:cat-O}
% =================================================================

\noindent
The natural conjecture (MC3): $D^b(\cO_{Y_\hbar(\mathfrak{sl}_N)})$
is thickly generated by finite-dimensional evaluation modules.

\begin{theorem}[The obstruction]
This is false.  Every object in
$\mathrm{thick}\langle\text{fd evals}\rangle$ has finite-dimensional
total cohomology.  Category $\cO$ contains Verma modules with
infinite-dimensional cohomology.  Therefore
$\mathrm{thick}\langle\text{fd evals}\rangle \neq D^b(\cO)$.
\end{theorem}

The failure is clean, immediate, and devastating.  It forced a
complete rethinking of MC3.

\subsection{The fork}

\textbf{Fork~A: Enlarge the generators.}  Add asymptotic modules,
prefundamental modules, or evaluation pullbacks of full BGG
$\cO(\fg)$.

\textbf{Fork~B: Enlarge the category.}  Move to
completed/coderived/ind-completed categories where
infinite-dimensional standards are totalizations of finite data.

The monograph takes Fork~A for type~$A$: the \emph{polynomial}
subcategory $\cO_{\mathrm{poly}}$ (modules with polynomial
character in weight variables) is proved to be thickly generated by
evaluation modules, via the Drinfeld classification.  The key
insight: modules outside $\cO_{\mathrm{poly}}$ have no quantum group
counterpart.  So factorization DK in its natural form is
unconditionally proved for type~$A$:
\[
\Fact_{\Eone}^{\mathrm{fd}}(Y_\hbar(\mathfrak{sl}_N))
\simeq
\Fact_{\Eone}^{\mathrm{fd}}(U_q(\mathfrak{sl}_N))^{\mathrm{op}}.
\]

The residual question: does $\cO = \cO_{\mathrm{poly}}$?

\subsection{The shifted-envelope dream}

The right ambient category may be the anti-dominantly shifted
envelope $\cO^{\mathrm{sh} \leq 0} = \bigoplus_{\mu \leq 0} \cO_\mu$,
where $\cO_\mu$ is Category~$\cO$ for the $\mu$-shifted Yangian
(Brundan--Kleshchev, Finkelberg--Tsymbaliuk).  Inside it, ordinary
$\cO$ lives at $\mu = 0$, negative prefundamental modules live at
$\mu = -\alpha_i^\vee$, and the Baxter $TQ$-relations should
categorify to exact triangles.  This is a fairy gate that has not
yet been opened.


% =================================================================
\section{The Fourier perspective}
\label{sec:fourier}
% =================================================================

\noindent
The bar construction is a Fourier transform.  This is not a metaphor.

The input is a chiral algebra (the ``function'').  The kernel is
$\eta_{ij} = d\log(z_i - z_j)$ (the ``exponential'').  The output
is a factorization coalgebra on $\Ran(X)$ (the ``spectrum'').
Verdier duality inverts the transform.  The cobar construction
$\Omega(\B(\cA))$ reconstructs $\cA$ from its spectrum.

\begin{theorem}[The four properties of the Fourier transform]
\label{thm:four-properties}
\begin{enumerate}[nosep, label=\textup{(\Alph*)}]
\item \textbf{Intertwining.}
  $\mathbb{D}_{\Ran} \circ \B_X \simeq \B_X \circ (-)^!$.
\item \textbf{Inversion.}
  On the Koszul locus, $\Omega(\B(\cA)) \xrightarrow{\sim} \cA$.
\item \textbf{Plancherel.}
  $Q_g(\cA) \oplus Q_g(\cA^!) \simeq
  R\Gamma(\overline{\cM}_g, \cZ_\cA)$.
\item \textbf{Characteristic function.}
  $\kappa(\cA)$ with $\hat{A}$-genus generating function.
\end{enumerate}
\end{theorem}

When $X = \Spec k$: configuration spaces collapse, the propagator
vanishes, $d_{\mathrm{dR}} = 0$, $d_{\mathrm{res}}$ reduces to the
binary bracket, and $\B_X(\cA)$ becomes
$T^c(s^{-1}\bar{\cA}, d_{\mathrm{bracket}})$---the classical bar
construction of Loday--Vallette.  The entire theory of operadic
Koszul duality is the zero-dimensional shadow of the chiral Fourier
transform.

\begin{remark}[The single deepest fact]
Verdier duality on $\overline{C}_n(X)$ exchanges residues at
collision divisors with inclusions of boundary strata.  The bar
differential (residues) and the cobar differential (inclusions) are
Verdier-dual operations.  Theorems~A--D are its consequences:
A~(intertwining) because Verdier duality interchanges
$\B(\cA)$ and $\B(\cA^!)$;
B~(inversion) because Verdier duality is an involution;
C~(complementarity) because the genus-$g$ obstructions of $\cA$ and
$\cA^!$ are Verdier-dual;
D~(characteristic) because the generating function of Verdier-dual
Chern classes is the $\hat{A}$-genus, by Grothendieck--Riemann--Roch.

Every computation in the monograph is an unwinding of this single
geometric principle in specific examples.
\end{remark}


% =================================================================
\section{The master conjecture hierarchy}
\label{sec:mc-hierarchy}
% =================================================================

\noindent
The ${\sim}100$ conjectures in the monograph reduce to five,
organized by logical dependency.

\medskip
\noindent
\textbf{MC1} (PBW concentration).
\emph{Proved} for KM, Virasoro, principal $\cW_N$.  The entry
theorem.  Without it, the interacting families remain ``conditionally
Koszul.''

\medskip
\noindent
\textbf{MC2} (cyclic $\Linf$ and $\Theta_\cA$).
\emph{Proved} (full resolution: chiral graph complex + modular-operadic
MC + tautological-line support).  The universal MC element exists,
is scalar-saturated, and has the right Verdier and clutching
compatibilities.

\medskip
\noindent
\textbf{MC3} (factorization DK/KL).
DK-0/1 proved.  DK-$1\frac{1}{2}$ (lattice) proved.
DK-2/3 type~$A$ proved.  Polynomial $\cO$ at all $N$ proved.
Category~$\cO$ at $N \geq 3$ open.  The lattice route bypasses
thick generation entirely via sectorwise finiteness.

\medskip
\noindent
\textbf{MC4} (infinite-tower comparison).
M-level complete for $\cW_\infty$ and Yangian RTT.  H-level
comparison open, reduced to coefficient identities:
$\mathsf{K}^{\mathrm{line}}_{a,b}(N) = \mathsf{K}^{\mathrm{RTT}}_{a,b}(N)$
(Yangian, detected on the boundary strip
$\{\Delta_{a,0}(N)\}_{0 \leq a \leq N}$) and
$\mathsf{C}^{\mathrm{res}}_{s,t;u;m,n}(N) =
\mathsf{C}^{\mathrm{DS}}_{s,t;u;m,n}(N)$
($\cW_\infty$, detected on the seed packet $\cI_N$).  At stage~$4$,
the live frontier is exactly six OPE coefficients organized as three
local blocks.

\medskip
\noindent
\textbf{MC5} (BV/BRST = bar).
Genus~$0$ proved.  Higher genus downstream.

\begin{remark}[Three phases of understanding]
Phase~1 (vision): ``Chiral Koszul duality lifts to curves.''  Broad
claims, few proofs.
Phase~2 (obstruction): literal MC3 dies, $d^2 \neq 0$ discovered,
enrichment error invalidates months of predictions.
Phase~3 (surgery): each MC reduced to explicit finite-detection
criteria.  The Yangian MC4 reduces to a residue on $e_1 \otimes e_2$
in $V \otimes V$.

The most ambitious theorem is the one stated with the most precision.
\end{remark}


% =================================================================
\section{Wisdom}
\label{sec:wisdom}
% =================================================================

\noindent
What follows is not formalized.  It is the residue of years of work,
extracted at the collision divisor of experience and ignorance.

\medskip
\noindent
\textit{1. The bar complex knows more than you tell it.}
The $\hat{A}$-genus, the Feigin--Frenkel involution, the
Kodaira--Spencer clutching compatibility---none are visible in the
OPE data.  They emerge from the geometry of $\overline{C}_n(X)$.
Geometric objects have richer structure than their algebraic
presentations suggest.  Trust the geometry.

\medskip
\noindent
\textit{2. Never decompose the differential by pole order.}
The split $d = d_{\mathrm{bracket}} + d_{\mathrm{curvature}}$ is
useful for understanding but fatal for computing.  The Borcherds
identity governs the full OPE simultaneously.  The Jacobi identity
governs only the simple pole.  When there is a double pole, these
are different identities.

\medskip
\noindent
\textit{3. The enrichment is $\fg \otimes \fg$, not $\Lambda^2(\fg)$.}
For $\mathfrak{sl}_2$: dimension $9$ versus dimension $3$.  Getting
this wrong makes every genus-$1$ prediction wrong by a factor that
grows with rank.  Several months were lost here.

\medskip
\noindent
\textit{4. ``$\Eone$-chiral'' $\neq$ ``topological $\Eone$.''}
An $\Eone$-chiral algebra lives on a curve (real dimension $2$).
A topological $\Eone$-algebra lives on an interval (real
dimension $1$).  The suffix ``-chiral'' is the difference between
the entire subject and nothing.

\medskip
\noindent
\textit{5. The curvature is the content, not the obstruction.}
$\dfib^2 = \kappa \cdot \omega_g$ is not a failure.  It is the
modular characteristic.  The anomaly \emph{is} the invariant.
$\kappa = 0$ is simultaneously uncurvedness and anomaly cancellation.

\medskip
\noindent
\textit{6. Free fields and interacting fields are different sports.}
The jump from one OPE pole to two is a phase transition.  Free-field
intuition does not transfer.  The Borcherds identity, the PBW
spectral sequence, the family-specific collapse proofs---the entire
machinery of the monograph exists to cross this gap.

\medskip
\noindent
\textit{7. Precision enables ambition.}
The discovery that the literal MC3 statement was false
\emph{strengthened} the programme by identifying the correct
statement ($\cO_{\mathrm{poly}}$ is the natural domain).  The
cubic Casimir error was corrected by a degree count, not a
calculation.  The most powerful theorems come from the most honest
accounting.

\medskip
\noindent
\textit{8. Status governance is mathematical practice.}
In a monograph with ${\sim}1400$ claims, bookkeeping errors are
indistinguishable from mathematical errors.  The concordance
(Chapter~34) is the constitution: when earlier chapters disagree,
they are subordinate.

\medskip
\noindent
\textit{9. The Fourier perspective unifies everything.}
The bar construction is a non-abelian integral transform.
Theorems A--D are its structural properties.  When $X$ degenerates
to a point, the transform degenerates to classical Koszul duality.
On an elliptic curve, the Heisenberg case literally reduces to
Fourier--Mukai.  The logarithmic propagator is the kernel; Verdier
duality is the inversion formula; the $\hat{A}$-genus is the
characteristic function.  This is not an analogy.  It is the
definition.

\medskip
\noindent
\textit{10. One geometric fact underlies everything.}
Verdier duality on $\overline{C}_n(X)$ exchanges residues at
collision divisors with inclusions of boundary strata.  Every
theorem, every computation, every identification in the monograph
is an unwinding of this single principle.

\bigskip
\begin{center}
$*\quad*\quad*$
\end{center}

\bigskip
\noindent\textit{The theorems are the stones.  The proofs are the
mortar.  But the building stands because of the invisible lines of
force between them---the insights that told us which stone to place
next, and which wall was load-bearing, and where the ground would
give way.  These notes record the invisible lines.}

\end{document}
